\chapter{Efeitos eletrônicos e intermediários reativos}\label{ch:b1:electronic-effects}

O ácido trifluoroetanoico, \ce{CF3COOH}, é um ácido cerca de dez mil
vezes mais forte que o ácido etanoico, \ce{CH3COOH}, embora a ligação
O--H ácida esteja a três ligações dos átomos de flúor. O fenol é um
milhão de vezes mais ácido que o etanol, e uma molécula de amônia ligada
a um anel benzênico --- a anilina --- é um milhão de vezes menos básica
que uma ligada a um grupo metila. A química orgânica está cheia de
contrastes assim, e dois efeitos explicam a maioria deles: a atração dos
átomos eletronegativos ao longo das ligações $\sigma$ e o espalhamento
dos pares de elétrons sobre os sistemas $\pi$. Os mesmos dois efeitos
decidem quais ligações se rompem, quais intermediários se formam e onde
os reagentes atacam: eles são a gramática dos \omterm{def:b1:elementary-steps:elementary-step}{mecanismos de reação} dos
próximos capítulos.

\begin{recall}
A \omterm{def:b1:periodicity:electronegativity}{eletronegatividade} e suas tendências (\cref{ch:b1:periodicity}); as
\omterm{def:b1:lewis-resonance-vsepr:lewis}{estruturas de Lewis}, as \omterm{def:b1:lewis-resonance-vsepr:formal-charge}{cargas formais} e a ressonância
(\cref{ch:b1:lewis-resonance-vsepr}); o $\mathrm{p}K_a$ e a força dos
ácidos e das bases (\cref{ch:b1:predominant-reaction}). O volume escolar
(ano 12) desenhou \omterm{def:b1:electronic-effects:curly-arrow}{setas curvas} de um \omterm{def:b1:electronic-effects:nucleophile}{nucleófilo} para um \omterm{def:b1:electronic-effects:nucleophile}{eletrófilo}.
\end{recall}

\section{Efeitos indutivos}

\begin{definition}[Efeito indutivo]\label{def:b1:electronic-effects:inductive}
O \emph{efeito indutivo}\index{efeito indutivo} é o deslocamento da
densidade eletrônica ao longo das ligações $\sigma$ em direção a um
átomo ou grupo eletronegativo. Um grupo atrator (F, Cl, O, \ce{NO2},
\ce{CF3}) tem efeito $-I$; os grupos alquila e os átomos com carga
negativa empurram elétrons, um efeito $+I$. O efeito se enfraquece
rapidamente com o número de ligações intermediárias.
\end{definition}

\begin{proposition}[Efeito indutivo e acidez]\label{prop:b1:electronic-effects:inductive-pka}
Um grupo atrator perto do sítio ácido de um ácido carboxílico estabiliza
o íon carboxilato espalhando sua carga negativa e diminui o
$\mathrm{p}K_a$: quanto mais grupos atratores, e quanto mais próximos,
mais forte é o ácido.
\end{proposition}

\begin{proof}
O $\mathrm{p}K_a$ mede a posição de $\mathrm{RCOOH} \rightleftharpoons
\mathrm{RCOO^-} + \ce{H+}$. Um grupo que retira densidade eletrônica do
carboxilato espalha sua carga por mais átomos, o que abaixa sua energia
mais que a do ácido neutro; o equilíbrio se desloca para a direita e
$K_a$ aumenta. A série medida o confirma.
\end{proof}

\begin{omfigure}
\begin{tikzpicture}
\begin{axis}[
  xbar, width=0.78\linewidth, height=5.2cm, xmin=0, xmax=5.4,
  symbolic y coords={TFA,TCA,DCA,MCA,AA}, ytick=data,
  yticklabels={\ce{CF3COOH}, \ce{CCl3COOH}, \ce{CHCl2COOH}, \ce{CH2ClCOOH}, \ce{CH3COOH}},
  xlabel={$\mathrm{p}K_a$ a \qty{25}{\celsius}}, bar width=9pt,
  nodes near coords, nodes near coords align={horizontal},
  every node near coord/.append style={font=\scriptsize},
  tick label style={font=\small}, label style={font=\small}, enlarge y limits=0.15]
\addplot[fill=omDef!45, draw=omDef] coordinates {(0.52,TFA) (0.51,TCA) (1.35,DCA) (2.87,MCA) (4.76,AA)};
\end{axis}
\end{tikzpicture}

{\small Os ácidos etanoicos halogenados. Cada cloro abaixa o
$\mathrm{p}K_a$; com três, o ácido é cerca de dez mil vezes mais forte
que o ácido etanoico. Os ácidos trifluoro- e tricloroetanoico têm a
mesma força, dentro da incerteza das medidas para ácidos tão fortes.}
% ledger: pka:ethanoic, pka:chloroethanoic, pka:dichloroethanoic, pka:trichloroethanoic, pka:trifluoroethanoic
\end{omfigure}

Os grupos alquila agem no sentido oposto, um pouco: o ácido propanoico
($\mathrm{p}K_a
= 4.89$) é ligeiramente mais fraco que o ácido etanoico (4.76), e o
ácido metanoico (3.74), sem grupo alquila, é mais forte que ambos.
% ledger: pka:propanoic, pka:methanoic

\section{Efeitos mesoméricos}

Quando um átomo que leva um \omterm{def:b1:lewis-resonance-vsepr:lewis}{par não ligante}, ou uma carga, está ligado a
um sistema $\pi$, seus elétrons podem ser compartilhados com esse
sistema; as \omterm{def:b1:lewis-resonance-vsepr:resonance}{estruturas de ressonância} (\cref{ch:b1:lewis-resonance-vsepr})
descrevem esse compartilhamento.

\begin{definition}[Efeito mesomérico, grupos doadores e aceptores]\label{def:b1:electronic-effects:mesomeric}
O \emph{efeito mesomérico}\index{efeito mesomerico@efeito mesomérico} é o deslocamento da
densidade eletrônica através de um sistema $\pi$ conjugado, mostrado
pelas \omterm{def:b1:lewis-resonance-vsepr:resonance}{estruturas de ressonância}. Um \emph{grupo doador}\index{grupo doador}
cede um \omterm{def:b1:lewis-resonance-vsepr:lewis}{par não ligante} ao sistema, um efeito $+M$ (\ce{-OH},
\ce{-O^-}, \ce{-NH2}, \ce{-OR}); um \emph{grupo aceptor}\index{grupo aceptor}
retira elétrons dele por meio de uma ligação $\pi$ com um átomo
eletronegativo, um efeito $-M$ (\ce{-NO2},
\ce{-C=O}, \ce{-C#N}).
\end{definition}

\begin{omfigure}
\setchemfig{atom sep=2em}
\schemestart
\chemfig{CH_3-C(=[1]O)-[7]O^{-}}
\arrow{<->}
\chemfig{CH_3-C(-[1]O^{-})=[7]O}
\schemestop
\hspace{2.5em}
\schemestart
\chemfig{*6(=-=(-O^{-})-=-)}
\arrow{<->}
\chemfig{*6((=O)-=-\chemabove{}{\scriptstyle\ominus}-=-)}
\schemestop

\medskip
\schemestart
\chemfig{H_3C-C(=[1]O)-[7]NH_2}
\arrow{<->}
\chemfig{H_3C-C(-[1]O^{-})=[7]NH_2^{+}}
\schemestop

{\small \omterm{def:b1:lewis-resonance-vsepr:resonance}{Estruturas de ressonância} do íon etanoato (a carga compartilhada
por dois oxigênios equivalentes), do íon fenóxido (a carga espalhada no
anel; uma das três estruturas do anel está representada) e de uma amida
(o \omterm{def:b1:lewis-resonance-vsepr:lewis}{par não ligante} do nitrogênio compartilhado com a C=O, o que faz das
amidas \omterm{def:b1:predominant-reaction:strong-acid}{bases fracas}).}
\end{omfigure}

\begin{proposition}[Efeito mesomérico e acidez]\label{prop:b1:electronic-effects:mesomeric-pka}
Um ácido cuja base conjugada espalha sua carga por ressonância é mais
forte que um ácido cuja base não pode fazê-lo: os ácidos carboxílicos
($\mathrm{p}K_a$ em torno de 4--5) são muito mais fortes que os álcoois
(cerca de 16), e o fenol (10.0) fica entre eles. Um \omterm{def:b1:electronic-effects:mesomeric}{grupo aceptor} na
posição \textit{orto} ou \textit{para} em relação ao OH de um fenol o
torna mais forte do que na posição \textit{meta}.
\end{proposition}

\begin{proof}
Em um alcóxido, a carga fica em um único oxigênio; em um carboxilato,
ela é compartilhada por dois oxigênios; em um fenóxido, ela é
compartilhada pelo oxigênio e por três carbonos do anel, menos
eletronegativos que o oxigênio, um ganho menor. Um grupo \ce{NO2} no
carbono \textit{orto} ou \textit{para} pode receber ele mesmo a carga
por meio de uma \omterm{def:b1:lewis-resonance-vsepr:resonance}{estrutura de ressonância} com \ce{C=N+(-O^-)-O^-}; em
\textit{meta}, nenhuma \omterm{def:b1:lewis-resonance-vsepr:resonance}{estrutura de ressonância} leva a carga ao carbono
que porta \ce{NO2}, e só resta o \omterm{def:b1:electronic-effects:inductive}{efeito indutivo}. Os valores medidos:
2-nitrofenol 7.23, 4-nitrofenol 7.16, 3-nitrofenol 8.36, contra 9.99
para o fenol e 15.9 para o etanol.
\end{proof}
% ledger: pka:ethanol, pka:phenol, pka:2-nitrophenol, pka:3-nitrophenol, pka:4-nitrophenol

O mesmo raciocínio se aplica às bases. A metilamina ($\mathrm{p}K_a$ de
seu amônio 10.66) é uma base mais forte que a amônia (9.25,
\cref{ch:b1:predominant-reaction}), com o grupo metila empurrando
elétrons para o nitrogênio; na anilina, o \omterm{def:b1:lewis-resonance-vsepr:lewis}{par não ligante} é
compartilhado com o anel e o amônio é muito mais ácido (4.60). A
piridina (5.23) conserva seu \omterm{def:b1:lewis-resonance-vsepr:lewis}{par não ligante}, mas em um nitrogênio preso
em um anel plano, com mais caráter \textit{s}, menos disponível que em
uma amina.
% ledger: pka:methylammonium, pka:anilinium, pka:pyridinium, dfg:NH4+_ao

\section{Romper uma ligação: os intermediários reativos}

\begin{definition}[Homólise e heterólise]\label{def:b1:electronic-effects:cleavage}
Uma \omterm{def:b1:lewis-resonance-vsepr:lewis}{ligação covalente} se rompe por \emph{homólise}\index{homolise@homólise} quando cada
átomo fica com um dos dois elétrons de ligação, dando dois \omterm{def:b1:electronic-effects:intermediates}{radicais}, e
por \emph{heterólise}\index{heterolise@heterólise} quando um átomo fica com os dois,
dando um cátion e um ânion.
\end{definition}

\begin{definition}[Intermediários reativos]\label{def:b1:electronic-effects:intermediates}
Um \emph{intermediário reativo}\index{intermediario reativo@intermediário reativo} é uma espécie
formada em uma etapa de um mecanismo e consumida em uma etapa posterior,
reativa demais para se acumular (\cref{ch:b1:elementary-steps}). Um
\emph{carbocátion}\index{carbocation@carbocátion} tem um carbono com três ligações e uma
carga positiva (seis \omterm{def:b1:quantum-numbers:configuration}{elétrons de valência}); um
\emph{carbânion}\index{carbanion@carbânion}, um carbono com três ligações, um \omterm{def:b1:lewis-resonance-vsepr:lewis}{par não
ligante} e uma carga negativa; um \emph{radical}\index{radical}, um átomo
com um \omterm{def:b1:quantum-numbers:unpaired}{elétron desemparelhado}.
\end{definition}

\begin{definition}[Classe de um carbono]\label{def:b1:electronic-effects:carbon-class}
Um carbono ligado a um único outro carbono é um \emph{carbono
primário}\index{carbono primario@carbono primário}, a dois, um \emph{carbono
secundário}\index{carbono secundario@carbono secundário}, a três, um \emph{carbono
terciário}\index{carbono terciario@carbono terciário}. Um \omterm{def:b1:electronic-effects:intermediates}{carbocátion}, um \omterm{def:b1:electronic-effects:intermediates}{radical} ou um
haloalcano recebe a classe do carbono considerado.
\end{definition}

\begin{omfigure}
\begin{tikzpicture}[font=\small]
  % carbocation: planar, empty p orbital
  \begin{scope}
    \fill[omProp!15, draw=omProp] (0,0.15) ellipse (0.18 and 0.75);
    \fill[omProp!15, draw=omProp] (0,-0.15) ellipse (0.18 and 0.75);
    \draw[thick] (0,0) -- (-1.0,-0.25) node[left] {R};
    \draw[thick] (0,0) -- (0.95,-0.35) node[right] {R};
    \draw[thick] (0,0) -- (0.35,0.45) node[above right] {R};
    \node[fill=white, inner sep=1pt] at (0,0) {\ce{C+}};
    \node at (0,-1.75) {carbocátion: plano,};
    \node at (0,-2.15) {orbital \textit{p} vazio};
  \end{scope}
  % radical: nearly planar, one electron
  \begin{scope}[xshift=4.6cm]
    \draw[omMeth] (0,0.15) ellipse (0.18 and 0.75);
    \draw[omMeth] (0,-0.15) ellipse (0.18 and 0.75);
    \fill (0,0.55) circle (0.05);
    \draw[thick] (0,0) -- (-1.0,-0.25) node[left] {R};
    \draw[thick] (0,0) -- (0.95,-0.35) node[right] {R};
    \draw[thick] (0,0) -- (0.35,0.45) node[above right] {R};
    \node[fill=white, inner sep=1pt] at (0,0) {C};
    \node at (0,-1.75) {radical: quase plano,};
    \node at (0,-2.15) {um elétron};
  \end{scope}
  % carbanion: pyramidal with a lone pair
  \begin{scope}[xshift=9.2cm]
    \draw[thick] (0,0) -- (-0.95,-0.55) node[left] {R};
    \draw[thick] (0,0) -- (0.95,-0.55) node[right] {R};
    \draw[thick] (0,0) -- (0.2,-0.9) node[below] {R};
    \fill[omDef!15, draw=omDef] (0,0.55) ellipse (0.22 and 0.45);
    \fill (-0.07,0.65) circle (0.045); \fill (0.07,0.65) circle (0.045);
    \node[fill=white, inner sep=1pt] at (0,0) {\ce{C-}};
    \node at (0,-1.75) {carbânion: piramidal,};
    \node at (0,-2.15) {par não ligante};
  \end{scope}
\end{tikzpicture}

{\small Geometrias dos três \omterm{def:b1:electronic-effects:intermediates}{intermediários reativos} do carbono (R: H ou
um grupo alquila). O \omterm{def:b1:electronic-effects:intermediates}{carbocátion} plano pode ser atacado por qualquer uma
das faces, um ponto que importa para a estereoquímica do
\cref{ch:b1:nucleophilic-substitution}.}
\end{omfigure}

\begin{proposition}[Estabilidade dos carbocátions]\label{prop:b1:electronic-effects:carbocation-order}
Os \omterm{def:b1:electronic-effects:intermediates}{carbocátions} são estabilizados pelos grupos que lhes cedem elétrons:
terciário $>$ secundário $>$ primário $>$ metila; um \omterm{def:b1:electronic-effects:intermediates}{carbocátion} vizinho
de uma ligação dupla (alílico) ou de um anel benzênico (benzílico) é
ainda estabilizado por ressonância, e um vizinho de um átomo com \omterm{def:b1:lewis-resonance-vsepr:lewis}{par não
ligante} (\ce{R-O-CH2+}), mais ainda.
\end{proposition}

\begin{proof}
Cada grupo alquila empurra densidade eletrônica para o carbono pobre em
elétrons (efeito $+I$, e uma interação aparentada de suas ligações C--H
com o orbital vazio, descrita no volume do 2.º~ano). Um cátion alílico
tem duas \omterm{def:b1:lewis-resonance-vsepr:resonance}{estruturas de ressonância} que põem a carga em um ou no outro
carbono da extremidade; um cátion benzílico tem quatro, três delas no
anel; um \omterm{def:b1:lewis-resonance-vsepr:lewis}{par não ligante} do oxigênio forma uma quarta ligação com o
carbono catiônico, \ce{R-O+=CH2}, na qual todo átomo tem um octeto.
\end{proof}

Os \omterm{def:b1:electronic-effects:intermediates}{carbânions} seguem a ordem oposta (metila $>$ primário $>$ secundário
$>$ terciário) e são estabilizados pelos \omterm{def:b1:electronic-effects:mesomeric}{grupos aceptores}; os \omterm{def:b1:electronic-effects:intermediates}{radicais}
seguem a mesma ordem que os \omterm{def:b1:electronic-effects:intermediates}{carbocátions}, de modo menos marcado.

\section{Setas curvas, nucleófilos e eletrófilos}

\begin{definition}[Setas curvas]\label{def:b1:electronic-effects:curly-arrow}
Em um mecanismo, uma \emph{seta curva}\index{seta curva} mostra o
deslocamento de um par de elétrons: ela parte de um \omterm{def:b1:lewis-resonance-vsepr:lewis}{par não ligante} ou de
uma ligação e termina em um átomo (onde se forma um \omterm{def:b1:lewis-resonance-vsepr:lewis}{par não ligante} ou
uma nova ligação) ou entre dois átomos (uma nova ligação). Uma
\emph{seta de meia ponta}\index{seta de meia ponta} mostra o deslocamento de
um único elétron, nas reações radicalares.
\end{definition}

\begin{method}[Desenhar setas curvas]\label{met:b1:electronic-effects:curly-arrows}
\begin{enumerate}
  \item Parta dos elétrons: um \omterm{def:b1:lewis-resonance-vsepr:lewis}{par não ligante}, uma ligação $\pi$ ou uma
        ligação $\sigma$, nunca de uma carga positiva ou de um átomo sem
        elétrons.
  \item Termine onde o par vai: uma nova ligação entre os dois átomos, ou
        um \omterm{def:b1:lewis-resonance-vsepr:lewis}{par não ligante} no átomo mais eletronegativo quando uma
        ligação se rompe.
  \item Conte as cargas depois de cada etapa: o átomo que cede um \omterm{def:b1:lewis-resonance-vsepr:lewis}{par
        não ligante} para uma nova ligação ganha $+1$; o átomo que recebe
        um par como \omterm{def:b1:lewis-resonance-vsepr:lewis}{par não ligante} ganha $-1$; a carga total se
        conserva.
  \item Nunca exceda um octeto em C, N, O: quando um par chega, outro
        precisa sair.
\end{enumerate}
\end{method}

\begin{omfigure}
\setchemfig{atom sep=2.2em}
\schemestart
\chemfig{H_3C-C(-[2]H)(-[6]H)-[@{b}]@{br}Br}
\arrow{->[heterólise]}[,1.5]
\chemfig{H_3C-C(-[2]H)(-[6]H)^{+}}\+\chemfig{Br^{-}}
\schemestop
\chemmove{\draw[->, shorten <=2pt, shorten >=2pt](b) .. controls +(90:6mm) and +(90:6mm) .. (br);}

\medskip
\schemestart
\chemfig{H@{o}O^{-}}\+\chemfig{@{h}H-@{oc}O-H}
\arrow{<=>}
\chemfig{H_2O}\+\chemfig{HO^{-}}
\schemestop
\chemmove{\draw[->, thick, shorten <=2pt, shorten >=2pt](o) .. controls +(-80:8mm) and +(-100:8mm) .. (h);}
\par\vspace{5mm}

{\small \omterm{def:b1:electronic-effects:curly-arrow}{Setas curvas}. Em cima: \omterm{def:b1:electronic-effects:cleavage}{heterólise} da ligação C--Br, com o par
indo para o bromo. Embaixo: uma transferência de próton da água para o
hidróxido, escrita como o ataque de um \omterm{def:b1:lewis-resonance-vsepr:lewis}{par não ligante} da base ao
hidrogênio (a ligação O--H da água se rompe então, com seu par ficando
no oxigênio; a segunda seta fica para o
\cref{exo:b1:electronic-effects:4}).}
\end{omfigure}

\begin{definition}[Nucleófilo, eletrófilo, grupo de saída]\label{def:b1:electronic-effects:nucleophile}
Um \emph{nucleófilo}\index{nucleofilo@nucleófilo} é uma espécie que cede um par de
elétrons a um átomo que não seja o hidrogênio para formar uma ligação
(uma \omterm{def:b1:lewis-resonance-vsepr:lewis-acid}{base de Lewis}); um \emph{eletrófilo}\index{eletrofilo@eletrófilo} é uma espécie
que aceita esse par (um \omterm{def:b1:lewis-resonance-vsepr:lewis-acid}{ácido de Lewis}, muitas vezes um carbono com uma
carga parcial positiva). A \emph{nucleofilicidade}\index{nucleofilicidade} de
uma espécie é sua reatividade como nucleófilo, medida por \omterm{def:b1:rate-laws:order}{constantes de
velocidade}. Um \emph{grupo de saída}\index{grupo de saida@grupo de saída} é o grupo que
parte com o par de ligação quando uma ligação com o carbono se rompe por
\omterm{def:b1:electronic-effects:cleavage}{heterólise}.
\end{definition}

\begin{proposition}[Nucleofilicidade e basicidade]\label{prop:b1:electronic-effects:nucleophilicity-basicity}
Para \omterm{def:b1:electronic-effects:nucleophile}{nucleófilos} que atacam pelo mesmo átomo, a \omterm{def:b1:electronic-effects:nucleophile}{nucleofilicidade}
acompanha a basicidade: \ce{CH3O-} $>$ \ce{OH-} $>$ \ce{CH3COO-} $>$
\ce{H2O}. Ela falha ao longo de uma coluna da tabela periódica em
\omterm{def:b1:intermolecular-forces-solvents:solvent-classes}{solventes próticos}, nos quais \omterm{def:b1:electronic-effects:nucleophile}{nucleófilos} grandes, polarizáveis e pouco
solvatados reagem mais depressa, embora sejam bases mais fracas
(\ce{I-} $>$ \ce{Br-} $>$ \ce{Cl-} $>$ \ce{F-}), e para bases volumosas,
que são maus \omterm{def:b1:electronic-effects:nucleophile}{nucleófilos}.
\end{proposition}

\begin{proof}
A basicidade mede o equilíbrio do ataque a um próton; a
\omterm{def:b1:electronic-effects:nucleophile}{nucleofilicidade}, a velocidade do ataque a um carbono; ambas aumentam
com a disponibilidade do par de elétrons, daí o paralelo para um mesmo
átomo. A velocidade também depende da energia necessária para retirar o
solvente do \omterm{def:b1:electronic-effects:nucleophile}{nucleófilo} e da deformação de sua nuvem eletrônica em
direção a um carbono congestionado: ânions pequenos como \ce{F-} estão
fortemente ligados às moléculas de \omterm{def:b1:intermolecular-forces-solvents:solvent-classes}{solvente prótico}
(\cref{ch:b1:intermolecular-forces-solvents}), e os grupos volumosos
dificultam a aproximação de um carbono, mas não de um próton.
\end{proof}

Um bom \omterm{def:b1:electronic-effects:nucleophile}{grupo de saída} é uma \omterm{def:b1:predominant-reaction:strong-acid}{base fraca}, estável depois de levar o par:
\ce{I-}, \ce{Br-}, \ce{Cl-}, a água, os sulfonatos; \ce{OH-} e \ce{RO-},
\omterm{def:b1:predominant-reaction:strong-acid}{bases fortes}, são maus \omterm{def:b1:electronic-effects:nucleophile}{grupos de saída} --- é por isso que os álcoois
precisam ser ativados antes de uma substituição
(\cref{ch:b1:alcohol-activation}).

\section{Exercícios}

\begin{exercise}[$\star$]\label{exo:b1:electronic-effects:1}
Classifique cada carbono do 2-metilbutano e do 2-bromo-2-metilpropano
como primário, secundário, terciário ou quaternário (ligado a quatro
carbonos). Que \omterm{def:b1:electronic-effects:intermediates}{carbocátion} se forma com mais facilidade a partir de cada
brometo de fórmula \ce{C4H9Br}?
\end{exercise}

\begin{exercise}[$\star$]\label{exo:b1:electronic-effects:2}
A partir dos valores de $\mathrm{p}K_a$, calcule a razão entre as
\omterm{def:b1:predominant-reaction:acidity-constant}{constantes de acidez} dos ácidos cloroetanoico e etanoico, e dos ácidos
trifluoroetanoico e etanoico.
\end{exercise}

\begin{exercise}[$\star$]\label{exo:b1:electronic-effects:3}
Desenhe as \omterm{def:b1:lewis-resonance-vsepr:resonance}{estruturas de ressonância} do íon etanoato, do íon
4-nitrofenóxido (com a carga em um oxigênio do grupo nitro) e do cátion
alila \ce{CH2=CH-CH2+}.
\end{exercise}

\begin{exercise}[$\star$]\label{exo:b1:electronic-effects:4}
Complete as \omterm{def:b1:electronic-effects:curly-arrow}{setas curvas} da transferência de próton da figura e de
\ce{NH3 + H3O+ -> NH4+ + H2O}. Verifique as cargas.
\end{exercise}

\begin{exercise}[$\star\star$]\label{exo:b1:electronic-effects:5}
Ordene por acidez crescente: etanol, fenol, ácido etanoico,
4-nitrofenol, ácido tricloroetanoico, água ($\mathrm{p}K_a$ do par
\ce{H2O}/\ce{OH-}, 14.0). Justifique cada passo com um argumento
indutivo ou mesomérico.
\end{exercise}

\begin{exercise}[$\star\star$]\label{exo:b1:electronic-effects:6}
Ordene por basicidade crescente: anilina, metilamina, amônia, piridina.
Explique a posição da anilina com uma \omterm{def:b1:lewis-resonance-vsepr:resonance}{estrutura de ressonância}.
\end{exercise}

\begin{exercise}[$\star\star$]\label{exo:b1:electronic-effects:7}
Por que o 3-nitrofenol é mais fraco que o 4-nitrofenol, mas ainda mais
forte que o fenol? Desenhe as \omterm{def:b1:lewis-resonance-vsepr:resonance}{estruturas de ressonância} que justificam
as duas afirmações.
\end{exercise}

\begin{exercise}[$\star\star$]\label{exo:b1:electronic-effects:8}
Ordene os \omterm{def:b1:electronic-effects:intermediates}{carbocátions} \ce{CH3+}, \ce{(CH3)3C+}, \ce{CH2=CH-CH2+},
\ce{CH3CH2+}, \ce{C6H5CH2+}, \ce{CH3OCH2+}, e justifique.
\end{exercise}

\begin{exercise}[$\star\star$]\label{exo:b1:electronic-effects:9}
Identifique o \omterm{def:b1:electronic-effects:nucleophile}{nucleófilo} e o \omterm{def:b1:electronic-effects:nucleophile}{eletrófilo} em: \ce{CH3Br + OH-};
\ce{H2O + H+}; o etanal com um íon cianeto; \ce{BF3 + NH3}. Desenhe as
\omterm{def:b1:electronic-effects:curly-arrow}{setas curvas}.
\end{exercise}

\begin{exercise}[$\star\star\star$]\label{exo:b1:electronic-effects:10}
Uma solução contém \qty{0.010}{mol/L} de ácido tricloroetanoico. A
fórmula do \omterm{def:b1:predominant-reaction:strong-acid}{ácido fraco} $\mathrm{pH} = \frac12(\mathrm{p}K_a - \log C)$ é
válida? Calcule corretamente o pH e o \omterm{def:b1:predominant-reaction:dissociation}{grau de dissociação}.
\end{exercise}

\begin{exercise}[$\star\star\star$]\label{exo:b1:electronic-effects:11}
Os $\mathrm{p}K_a$ dos álcoois medidos em água aumentam do metanol
(15.5) para o etanol (15.9), o propan-2-ol (17.1) e o 2-metilpropan-2-ol
(19.2). Dê duas explicações, uma pelo \omterm{def:b1:electronic-effects:inductive}{efeito indutivo} dos grupos alquila
sobre o alcóxido, outra pela \omterm{def:b1:intermolecular-forces-solvents:dissolution}{solvatação} do alcóxido.
% ledger: pka:methanol, pka:ethanol, pka:2-propanol, pka:tBuOH
\end{exercise}

\begin{exercise}[$\star\star\star$]\label{exo:b1:electronic-effects:12}
Explique por que as amidas não são básicas no nitrogênio nem
nucleofílicas, enquanto as aminas são ambas as coisas, usando a
\omterm{def:b1:lewis-resonance-vsepr:resonance}{estrutura de ressonância} da amida. Que átomo de uma amida é protonado em
\omterm{def:b1:predominant-reaction:strong-acid}{ácido forte}?
\end{exercise}

\section{Problema: Por que o ácido trifluoroetanoico é tão forte}

\begin{problem}[{Problema de fim de semana --- a série indutiva dos ácidos
halogenados, os carboxilatos contra os alcóxidos, os três nitrofenóis, um
cálculo pela reação predominante e a razão entre as constantes de acidez
dos ácidos trifluoroetanoico e etanoico}]
\label{pb:b1:electronic-effects:1}
Dados ($\mathrm{p}K_a$ a \qty{25}{\celsius}): ácido etanoico 4.76,
cloroetanoico 2.87, dicloroetanoico 1.35, tricloroetanoico 0.51,
trifluoroetanoico 0.52, etanol 15.9, fenol 9.99, 2-nitrofenol 7.23,
3-nitrofenol 8.36, 4-nitrofenol 7.16; $\mathrm{p}K_e = 14.00$.

\medskip
\noindent\textbf{Parte I --- A série indutiva.}
\begin{enumerate}
  \item Calcule a razão entre os $K_a$ em cada passo, do ácido etanoico
        ao tricloroetanoico.
  \item O efeito de cada cloro é o mesmo? Comente.
  \item Explique a tendência com o \omterm{def:b1:electronic-effects:inductive}{efeito indutivo}.
  \item Por que um cloro no carbono de uma cadeia mais longa mais
        afastado de COOH quase não teria efeito?
  \item O flúor é mais eletronegativo que o cloro. O que mostra a
        comparação dos ácidos trifluoro- e tricloro-, e por que é difícil
        medir esses $\mathrm{p}K_a$ com precisão?
  \item Em que forma está cada um desses ácidos em pH 3?
\end{enumerate}

\medskip
\noindent\textbf{Parte II --- Carboxilatos e alcóxidos.}
\begin{enumerate}[resume]
  \item Desenhe as duas \omterm{def:b1:lewis-resonance-vsepr:resonance}{estruturas de ressonância} do íon etanoato.
  \item O que elas preveem para os dois comprimentos de ligação C--O do
        íon?
  \item Por que o íon etóxido não tem estruturas equivalentes?
  \item Calcule a razão entre os $K_a$ do ácido etanoico e do etanol.
  \item Situe o fenol entre os dois e justifique com as \omterm{def:b1:lewis-resonance-vsepr:resonance}{estruturas de
        ressonância} do fenóxido.
  \item Quais dentre o etanol, o fenol e o ácido etanoico reagem quase
        totalmente com uma solução de hidrogenocarbonato ($\mathrm{p}K_a$
        6.37 para
        \ce{CO2,H2O}/\ce{HCO3-})?
\end{enumerate}

\medskip
\noindent\textbf{Parte III --- Os nitrofenóis.}
\begin{enumerate}[resume]
  \item Ordene os três nitrofenóis e o fenol por acidez.
  \item Desenhe uma \omterm{def:b1:lewis-resonance-vsepr:resonance}{estrutura de ressonância} do 4-nitrofenóxido com a
        carga no grupo nitro.
  \item Mostre que não existe estrutura desse tipo para o
        3-nitrofenóxido.
  \item Por que o 3-nitrofenol ainda é mais ácido que o fenol?
  \item O 2-nitrofenol é ligeiramente mais fraco que o 4-nitrofenol,
        embora ambos tenham o \omterm{def:b1:electronic-effects:mesomeric}{efeito mesomérico}; proponha uma explicação
        que envolva seu OH e o grupo nitro vizinho.
  \item Em pH 7.5, quais nitrofenóis estão majoritariamente na forma
        aniônica? Por que isso faz do 4-nitrofenol um \omterm{def:b1:titration-methods:indicator}{indicador colorido}?
\end{enumerate}

\medskip
\noindent\textbf{Parte IV --- O ácido trifluoroetanoico na água.}
\begin{enumerate}[resume]
  \item Calcule o pH do ácido trifluoroetanoico a \qty{0.010}{mol/L}
        (não suponha dissociação fraca).
  \item Calcule o pH do ácido etanoico a \qty{0.010}{mol/L}.
  \item Escreva a reação do ácido trifluoroetanoico com o etanoato de
        sódio e calcule sua constante.
  \item Dê a razão $K_a(\ce{CF3COOH})/K_a(\ce{CH3COOH})$, como potência
        de dez e como número.
\end{enumerate}
\end{problem}
